\begin{definitionbox}{Definition --- Metric}
\begin{definition}[Metric]\label{def:metric-on-set}
Let \(X\) be a set.  A \emph{metric} on \(X\) is a real-valued function
\[
  d : X \times X \to \mathbb{R}
\]
such that for all \(x,y,z\in X\):
\begin{enumerate}
  \item \(d(x,y)\geq 0\), and \(d(x,y)=0\) if and only if \(x=y\);
  \item \(d(x,y)=d(y,x)\);
  \item \(d(x,z)\leq d(x,y)+d(y,z)\).
\end{enumerate}
\end{definition}
\end{definitionbox}

\begin{remark*}[Standard quantified statement]
\[
\begin{aligned}
&d:X\times X\to\mathbb{R}
\\
&{}\land
\forall x,y,z\in X{:}\;
d(x,y)\geq 0
\land
\bigl(d(x,y)=0\Longleftrightarrow x=y\bigr)
\\
&{}\land
d(x,y)=d(y,x)
\land
d(x,z)\leq d(x,y)+d(y,z).
\end{aligned}
\]
\end{remark*}

\begin{remark*}[Predicate reading]
\[
d(x,y)
\]
means the distance assigned by \(d\) from \(x\) to \(y\).
\end{remark*}

\begin{remark*}[Interpretation]
A metric assigns distances between points of a set.  Its axioms say that
distance is nonnegative, that zero distance identifies exactly equal points,
that distance is symmetric, and that the direct distance from \(x\) to \(z\)
is no larger than the distance obtained by passing through \(y\).
\end{remark*}

\begin{remark*}[Historical note]
This definition is modeled on Definition~1.1.1 of
{\'O} Searc{\'o}id's \emph{Metric Spaces} \cite{OSearcoidMetricSpaces}.
\end{remark*}

\LeanFormalizes{def:metric-on-set}{lra-lean}{LRA.VolumeIV.MetricSpaces.Foundations.MetricSpace}{MetricDefinition}{statement}

\DefinitionalRoot

\begin{definitionbox}{Definition --- Euclidean Distance on the Real Line}
\begin{definition}[Euclidean Distance on \(\mathbb{R}\)]
\label{def:real-line-euclidean-distance}
The \emph{Euclidean distance} on \(\mathbb{R}\) is the function
\[
  d_{\mathbb{R}}:\mathbb{R}\times\mathbb{R}\to\mathbb{R}_{\geq 0},
  \qquad
  d_{\mathbb{R}}(x,y)=|x-y|.
\]
\end{definition}
\end{definitionbox}

\begin{remark*}[Standard quantified statement]
\[
d_{\mathbb{R}}:\mathbb{R}\times\mathbb{R}\to\mathbb{R}_{\geq 0},
\qquad
\forall x,y\in\mathbb{R}{:}\quad d_{\mathbb{R}}(x,y)=|x-y|.
\]
\end{remark*}

\begin{remark*}[Predicate reading]
\[
\operatorname{IsMetric}(\mathbb{R},d_{\mathbb{R}})
\]
means that absolute difference is a metric on the real line.
\end{remark*}

\begin{remark*}[Interpretation]
The real-line distance records absolute displacement between two real numbers.
It forgets direction and keeps only separation.
\end{remark*}

\LeanFormalizes{def:real-line-euclidean-distance}{lra-lean}{LRA.VolumeIV.MetricSpaces.Metrics}{EuclideanRMetric}{statement}

\DefinitionalRoot

\begin{propositionbox}{Proposition}
\begin{proposition}[The Real-Line Euclidean Distance is a Metric]
\label{prop:real-line-euclidean-distance-metric}
The function \(d_{\mathbb{R}}(x,y)=|x-y|\) is a metric on \(\mathbb{R}\).

\smallskip
\noindent
\hyperref[prf:real-line-euclidean-distance-metric]{\textit{Go to proof.}}
\end{proposition}
\end{propositionbox}

\begin{remark*}[Standard quantified statement]
\[
  \operatorname{IsMetric}(\mathbb{R},d_{\mathbb{R}}).
\]
\end{remark*}

\begin{remark*}[Predicate reading]
\[
\operatorname{IsMetric}(\mathbb{R},d_{\mathbb{R}})
\]
means that absolute difference is a metric on the real line.
\end{remark*}

\begin{remark*}[Interpretation]
This proposition says that the usual absolute-value distance satisfies the
abstract metric axioms.
\end{remark*}

\LeanFormalizes{prop:real-line-euclidean-distance-metric}{lra-lean}{LRA.VolumeIV.MetricSpaces}{EuclideanDistanceIsAMetric}{checked}

\begin{dependencies}
\begin{itemize}
  \item \hyperref[def:real-line-euclidean-distance]{Definition: Euclidean Distance on \(\mathbb{R}\)}
  \item \hyperref[def:metric-space]{Definition: Metric Space}
\end{itemize}
\end{dependencies}

\begin{definitionbox}{Definition --- Euclidean Space}
\begin{definition}[Euclidean Space]\label{def:euclidean-space}
For \(n\in\mathbb{N}\), the \emph{\(n\)-dimensional Euclidean space} is
\(\mathbb{R}^n\) equipped with its standard coordinate structure, dot product,
Euclidean norm, and Euclidean distance.
\end{definition}
\end{definitionbox}

\begin{remark*}[Standard quantified statement]
For each \(n\in\mathbb{N}\), \(\operatorname{IsEuclideanSpace}(n)\) denotes
\(\mathbb{R}^n\) together with its standard coordinate structure, dot product,
Euclidean norm, and Euclidean distance.
\end{remark*}

\begin{remark*}[Predicate reading]
\[
\operatorname{IsEuclideanSpace}(n)
\]
means that \(\mathbb{R}^n\) carries its standard Euclidean coordinate,
inner-product, norm, and metric structure.
\end{remark*}

\begin{remark*}[Interpretation]
Euclidean space is the coordinate model of ordinary distance geometry.
\end{remark*}

\begin{remark*}[Historical note]
This convention is the setting for the Euclidean examples in Kainth's
\emph{A Comprehensive Textbook on Metric Spaces}
\cite{KainthComprehensiveMetricSpaces2023}.
\end{remark*}

\DefinitionalRoot

\begin{definitionbox}{Definition --- Euclidean Distance on \(\mathbb{R}^n\)}
\begin{definition}[Euclidean Distance on \(\mathbb{R}^n\)]
\label{def:euclidean-distance-rn}
For \(x=(x_1,\ldots,x_n)\) and \(y=(y_1,\ldots,y_n)\) in \(\mathbb{R}^n\),
the \emph{Euclidean distance} is the function
\[
  d_2:\mathbb{R}^n\times\mathbb{R}^n\to\mathbb{R}_{\geq 0},
  \qquad
  d_2(x,y)
  =
  \sqrt{(x_1-y_1)^2+\cdots+(x_n-y_n)^2}.
\]
\end{definition}
\end{definitionbox}

\begin{remark*}[Standard quantified statement]
\[
d_2:\mathbb{R}^n\times\mathbb{R}^n\to\mathbb{R}_{\geq 0},
\qquad
d_2(x,y)=\sqrt{\sum_{i=1}^n (x_i-y_i)^2}.
\]
\end{remark*}

\begin{remark*}[Interpretation]
The Euclidean distance measures the length of the coordinate difference vector
\(x-y\).
\end{remark*}

\DefinitionalRoot

\begin{theorembox}{Theorem}
\begin{theorem}[Cauchy--Schwarz Inequality]
\label{thm:cauchy-schwarz-inequality}
For every \(x,y\in\mathbb{R}^n\),
\[
  |x\cdot y|\leq |x|\,|y|.
\]
Equality holds if and only if \(x\) and \(y\) are linearly dependent.

\smallskip
\noindent
\hyperref[prf:cauchy-schwarz-inequality]{\textit{Go to proof.}}
\end{theorem}
\end{theorembox}

\begin{remark*}[Standard quantified statement]
\[
\forall x,y\in\mathbb{R}^n{:}\quad |x\cdot y|\leq |x|\,|y|.
\]
\end{remark*}

\begin{remark*}[Predicate reading]
\[
x,y\in\mathbb{R}^n
\Longrightarrow
|x\cdot y|\leq |x|\,|y|
\]
means that the absolute value of the dot product is bounded by the product of
Euclidean lengths.
\end{remark*}

\begin{remark*}[Interpretation]
Cauchy--Schwarz controls the dot product by the lengths of the two vectors.
It is the algebraic estimate behind Euclidean triangle inequalities.
\end{remark*}

\begin{dependencies}
\begin{itemize}
  \item \hyperref[def:euclidean-space]{Definition: Euclidean Space}
  \item \hyperref[def:euclidean-distance-rn]{Definition: Euclidean Distance on \(\mathbb{R}^n\)}
\end{itemize}
\end{dependencies}

\begin{corollarybox}{Corollary}
\begin{corollary}[Minkowski's Inequality]\label{cor:minkowski-inequality-rn}
For every \(x,y\in\mathbb{R}^n\),
\[
  |x+y|\leq |x|+|y|.
\]

\smallskip
\noindent
\hyperref[prf:minkowski-inequality-rn]{\textit{Go to proof.}}
\end{corollary}
\end{corollarybox}

\begin{remark*}[Standard quantified statement]
\[
\forall x,y\in\mathbb{R}^n{:}\quad |x+y|\leq |x|+|y|.
\]
\end{remark*}

\begin{remark*}[Predicate reading]
\[
x,y\in\mathbb{R}^n
\Longrightarrow
|x+y|\leq |x|+|y|
\]
means that the Euclidean length of a sum is at most the sum of the Euclidean
lengths.
\end{remark*}

\begin{remark*}[Interpretation]
Minkowski's inequality is the triangle inequality for Euclidean length.
\end{remark*}

\begin{dependencies}
\begin{itemize}
  \item \hyperref[thm:cauchy-schwarz-inequality]{Theorem: Cauchy--Schwarz Inequality}
  \item \hyperref[def:euclidean-distance-rn]{Definition: Euclidean Distance on \(\mathbb{R}^n\)}
\end{itemize}
\end{dependencies}

\begin{propositionbox}{Proposition}
\begin{proposition}[The Euclidean Distance on \(\mathbb{R}^n\) is a Metric]
\label{prop:euclidean-distance-rn-metric}
The function \(d_2\) is a metric on \(\mathbb{R}^n\).

\smallskip
\noindent
\hyperref[prf:euclidean-distance-rn-metric]{\textit{Go to proof.}}
\end{proposition}
\end{propositionbox}

\begin{remark*}[Standard quantified statement]
\[
  \operatorname{IsMetric}(\mathbb{R}^n,d_2).
\]
\end{remark*}

\begin{remark*}[Predicate reading]
\[
\operatorname{IsMetric}(\mathbb{R}^n,d_2)
\]
means that the Euclidean distance satisfies the metric axioms on
\(\mathbb{R}^n\).
\end{remark*}

\begin{remark*}[Interpretation]
The Euclidean distance makes \(\mathbb{R}^n\) into the standard finite
dimensional metric space.
\end{remark*}

\begin{dependencies}
\begin{itemize}
  \item \hyperref[def:metric-space]{Definition: Metric Space}
  \item \hyperref[def:euclidean-distance-rn]{Definition: Euclidean Distance on \(\mathbb{R}^n\)}
  \item \hyperref[cor:minkowski-inequality-rn]{Corollary: Minkowski's Inequality}
\end{itemize}
\end{dependencies}

\begin{definitionbox}{Definition --- Taxicab Metric}
\begin{definition}[Taxicab Metric]\label{def:taxicab-metric-rn}
For \(x,y\in\mathbb{R}^n\), the \emph{taxicab metric} is the function
\[
  d_1(x,y):=\sum_{i=1}^{n}|x_i-y_i|.
\]
\end{definition}
\end{definitionbox}

\begin{remark*}[Standard quantified statement]
\[
d_1:\mathbb{R}^n\times\mathbb{R}^n\to\mathbb{R}_{\geq 0},
\qquad
d_1(x,y)=\sum_{i=1}^{n}|x_i-y_i|.
\]
\end{remark*}

\begin{remark*}[Interpretation]
Taxicab distance measures coordinate displacement additively, as if travel
were constrained to coordinate-parallel streets.
\end{remark*}

\DefinitionalRoot

\begin{propositionbox}{Proposition}
\begin{proposition}[Taxicab Metric]\label{prop:taxicab-metric-rn}
The function \(d_1\) is a metric on \(\mathbb{R}^n\).

\smallskip
\noindent
\hyperref[prf:taxicab-metric-rn]{\textit{Go to proof.}}
\end{proposition}
\end{propositionbox}

\begin{remark*}[Standard quantified statement]
\[
\operatorname{IsMetric}(\mathbb{R}^n,d_1).
\]
\end{remark*}

\begin{remark*}[Predicate reading]
\[
\operatorname{IsMetric}(\mathbb{R}^n,d_1)
\]
means that the taxicab distance satisfies the metric axioms on
\(\mathbb{R}^n\).
\end{remark*}

\begin{remark*}[Interpretation]
The taxicab metric satisfies the metric axioms coordinate by coordinate, with
the real absolute-value triangle inequality summed over all coordinates.
\end{remark*}

\begin{dependencies}
\begin{itemize}
  \item \hyperref[def:metric-space]{Definition: Metric Space}
  \item \hyperref[def:taxicab-metric-rn]{Definition: Taxicab Metric}
\end{itemize}
\end{dependencies}

\begin{definitionbox}{Definition --- \(p\)-Metric on \(\mathbb{R}^n\)}
\begin{definition}[\(p\)-Metric on \(\mathbb{R}^n\)]\label{def:lp-metric-rn}
Let \(1\leq p<\infty\).  For \(x,y\in\mathbb{R}^n\), define
\[
  d_p(x,y)
  :=
  \left(\sum_{i=1}^{n}|x_i-y_i|^p\right)^{1/p}.
\]
\end{definition}
\end{definitionbox}

\begin{remark*}[Standard quantified statement]
\[
1\leq p<\infty
\Longrightarrow
d_p:\mathbb{R}^n\times\mathbb{R}^n\to\mathbb{R}_{\geq 0}.
\]
\end{remark*}

\begin{remark*}[Interpretation]
The \(p\)-metrics interpolate between taxicab distance at \(p=1\) and
Euclidean distance at \(p=2\), and they form the standard finite-dimensional
\(\ell^p\) family.
\end{remark*}

\DefinitionalRoot

\begin{theorembox}{Theorem}
\begin{theorem}[\(d_p\) is a Metric]\label{thm:lp-metric-rn}
For every \(1\leq p<\infty\), the function \(d_p\) is a metric on
\(\mathbb{R}^n\).

\smallskip
\noindent
\hyperref[prf:lp-metric-rn]{\textit{Go to proof.}}
\end{theorem}
\end{theorembox}

\begin{remark*}[Standard quantified statement]
\[
1\leq p<\infty
\Longrightarrow
\operatorname{IsMetric}(\mathbb{R}^n,d_p).
\]
\end{remark*}

\begin{remark*}[Predicate reading]
\[
\operatorname{IsMetric}(\mathbb{R}^n,d_p)
\]
means that the \(p\)-distance satisfies the metric axioms on \(\mathbb{R}^n\)
for \(1\leq p<\infty\).
\end{remark*}

\begin{remark*}[Interpretation]
The proof is the \(p\)-norm version of Minkowski's inequality.  The theorem
packages a whole family of metric structures on the same underlying set.
\end{remark*}

\begin{dependencies}
\begin{itemize}
  \item \hyperref[def:metric-space]{Definition: Metric Space}
  \item \hyperref[def:lp-metric-rn]{Definition: \(p\)-Metric on \(\mathbb{R}^n\)}
\end{itemize}
\end{dependencies}

\begin{definitionbox}{Definition --- Complex Metric}
\begin{definition}[Complex Metric]\label{def:complex-metric}
Identify \(\mathbb{C}\) with \(\mathbb{R}^2\).  The \emph{complex metric} is
\[
  d_{\mathbb{C}}(z,w):=|z-w|.
\]
\end{definition}
\end{definitionbox}

\begin{remark*}[Standard quantified statement]
\[
d_{\mathbb{C}}:\mathbb{C}\times\mathbb{C}\to\mathbb{R}_{\geq 0},
\qquad
d_{\mathbb{C}}(z,w)=|z-w|.
\]
\end{remark*}

\begin{remark*}[Interpretation]
As a distance function, the complex metric is the Euclidean metric on
\(\mathbb{R}^2\) written in complex notation.
\end{remark*}

\LeanFormalizes{def:complex-metric}{lra-lean}{LRA.VolumeIV.MetricSpaces.Metrics}{ComplexModulusMetric}{statement}

\DefinitionalRoot

\begin{propositionbox}{Proposition}
\begin{proposition}[The Complex Modulus Distance is a Metric]
\label{prop:complex-modulus-metric}
The function \(d_{\mathbb{C}}(z,w)=|z-w|\) is a metric on \(\mathbb{C}\).

\smallskip
\noindent
\hyperref[prf:complex-modulus-metric]{\textit{Go to proof.}}
\end{proposition}
\end{propositionbox}

\begin{remark*}[Standard quantified statement]
\[
  \operatorname{IsMetric}(\mathbb{C},d_{\mathbb{C}}).
\]
\end{remark*}

\begin{remark*}[Predicate reading]
\[
\operatorname{IsMetric}(\mathbb{C},d_{\mathbb{C}})
\]
means that modulus distance satisfies the metric axioms on the complex plane.
\end{remark*}

\begin{remark*}[Interpretation]
The modulus metric is the Euclidean metric of the complex plane written in
complex notation.
\end{remark*}

\LeanFormalizes{prop:complex-modulus-metric}{lra-lean}{LRA.VolumeIV.MetricSpaces}{ModulusIsAMetricOnTheComplexNumbers}{checked}

\begin{dependencies}
\begin{itemize}
  \item \hyperref[def:complex-metric]{Definition: Complex Metric}
  \item \hyperref[def:metric-space]{Definition: Metric Space}
\end{itemize}
\end{dependencies}

\begin{definitionbox}{Definition --- Real Circle Chord Metric}
\begin{definition}[Real Circle Chord Metric]\label{def:real-circle-chord-metric}
Let
\[
  S^1_{\mathbb{R}}
  :=
  \{(x,y)\in\mathbb{R}^2 : x^2+y^2=1\}.
\]
For \(p,q\in S^1_{\mathbb{R}}\), the \emph{real circle chord metric} is
\[
  d_{\mathrm{chord},\mathbb{R}}(p,q):=d_2(p,q),
\]
the Euclidean distance between the two points as points of \(\mathbb{R}^2\).
\end{definition}
\end{definitionbox}

\begin{remark*}[Standard quantified statement]
\[
d_{\mathrm{chord},\mathbb{R}}:
S^1_{\mathbb{R}}\times S^1_{\mathbb{R}}\to\mathbb{R}_{\geq 0},
\qquad
d_{\mathrm{chord},\mathbb{R}}(p,q)=d_2(p,q).
\]
\end{remark*}

\begin{remark*}[Interpretation]
Chord distance measures separation along the straight line segment joining two
points of the unit circle, not along the circular arc between them.
\end{remark*}

\LeanFormalizes{def:real-circle-chord-metric}{lra-lean}{LRA.VolumeIV.MetricSpaces.Metrics}{RealCircleChordMetric}{statement}

\DefinitionalRoot

\begin{propositionbox}{Proposition}
\begin{proposition}[The Real Circle Chord Distance is a Metric]
\label{prop:real-circle-chord-metric}
The function \(d_{\mathrm{chord},\mathbb{R}}(p,q)=d_2(p,q)\) is a metric on
\(S^1_{\mathbb{R}}\).

\smallskip
\noindent
\hyperref[prf:real-circle-chord-metric]{\textit{Go to proof.}}
\end{proposition}
\end{propositionbox}

\begin{remark*}[Standard quantified statement]
\[
  \operatorname{IsMetric}(S^1_{\mathbb{R}},d_{\mathrm{chord},\mathbb{R}}).
\]
\end{remark*}

\begin{remark*}[Interpretation]
The real chord metric is the Euclidean metric on \(\mathbb{R}^2\) restricted
to the unit circle.
\end{remark*}

\begin{dependencies}
\begin{itemize}
  \item \hyperref[def:euclidean-distance-rn]{Definition: Euclidean Distance on \(\mathbb{R}^n\)}
  \item \hyperref[def:real-circle-chord-metric]{Definition: Real Circle Chord Metric}
  \item \hyperref[def:metric-space]{Definition: Metric Space}
  \item \hyperref[prop:euclidean-distance-rn-metric]{Proposition: The Euclidean Distance on \(\mathbb{R}^n\) is a Metric}
\end{itemize}
\end{dependencies}

\begin{definitionbox}{Definition --- Complex Circle Chord Metric}
\begin{definition}[Complex Circle Chord Metric]\label{def:complex-circle-chord-metric}
Let
\[
  S^1_{\mathbb{C}}
  :=
  \{z\in\mathbb{C}: |z|=1\}.
\]
For \(z,w\in S^1_{\mathbb{C}}\), the \emph{complex circle chord metric} is
\[
  d_{\mathrm{chord},\mathbb{C}}(z,w):=|z-w|.
\]
\end{definition}
\end{definitionbox}

\begin{remark*}[Standard quantified statement]
\[
d_{\mathrm{chord},\mathbb{C}}:
S^1_{\mathbb{C}}\times S^1_{\mathbb{C}}\to\mathbb{R}_{\geq 0},
\qquad
d_{\mathrm{chord},\mathbb{C}}(z,w)=|z-w|.
\]
\end{remark*}

\begin{remark*}[Interpretation]
The complex circle chord metric is the complex-plane modulus distance
restricted to the unit circle.
\end{remark*}

\LeanFormalizes{def:complex-circle-chord-metric}{lra-lean}{LRA.VolumeIV.MetricSpaces.Metrics}{ComplexCircleChordMetric}{statement}

\DefinitionalRoot

\begin{propositionbox}{Proposition}
\begin{proposition}[The Complex Circle Admits the Chord Metric]
\label{prop:complex-circle-chord-metric}
The function \(d_{\mathrm{chord},\mathbb{C}}(z,w)=|z-w|\) is a metric
on \(S^1_{\mathbb{C}}\).

\smallskip
\noindent
\hyperref[prf:complex-circle-chord-metric]{\textit{Go to proof.}}
\end{proposition}
\end{propositionbox}

\begin{remark*}[Standard quantified statement]
\[
  \operatorname{IsMetric}(S^1_{\mathbb{C}},d_{\mathrm{chord},\mathbb{C}}).
\]
\end{remark*}

\begin{remark*}[Predicate reading]
\[
\operatorname{IsMetric}(S^1_{\mathbb{C}},d_{\mathrm{chord},\mathbb{C}})
\]
means that chord distance satisfies the metric axioms on the complex unit
circle.
\end{remark*}

\begin{remark*}[Interpretation]
The chord metric is obtained by restricting the complex modulus metric to the
unit circle.  Restricting a metric to a subset preserves the metric axioms.
\end{remark*}

\LeanFormalizes{prop:complex-circle-chord-metric}{lra-lean}{LRA.VolumeIV.MetricSpaces}{CircleAdmitsChordAsMetric}{checked}

\begin{dependencies}
\begin{itemize}
  \item \hyperref[def:complex-metric]{Definition: Complex Metric}
  \item \hyperref[def:complex-circle-chord-metric]{Definition: Complex Circle Chord Metric}
  \item \hyperref[def:metric-space]{Definition: Metric Space}
  \item \hyperref[prop:complex-plane-metric-space]{Proposition: The Complex Plane is a Metric Space}
\end{itemize}
\end{dependencies}

\begin{definitionbox}{Definition --- Supremum Metric on Bounded Real Sequences}
\begin{definition}[Supremum Metric on Bounded Real Sequences]
\label{def:bounded-real-sequence-space}
Let \(\ell^\infty(\mathbb{R})\) denote the set of all bounded real sequences:
\[
  \ell^\infty(\mathbb{R})
  :=
  \left\{x=(x_n)_{n\in\mathbb{N}}:
  \exists M\geq 0\;\forall n\in\mathbb{N}\; |x_n|\leq M\right\}.
\]
For \(x,y\in\ell^\infty(\mathbb{R})\), define
\[
  d_\infty(x,y)
  :=
  \sup_{n\in\mathbb{N}} |x_n-y_n|.
\]
\end{definition}
\end{definitionbox}

\begin{remark*}[Standard quantified statement]
\[
d_\infty:
\ell^\infty(\mathbb{R})\times\ell^\infty(\mathbb{R})
\to
\mathbb{R}_{\geq 0},
\qquad
d_\infty(x,y)=\sup_{n\in\mathbb{N}} |x_n-y_n|.
\]
\end{remark*}

\begin{remark*}[Interpretation]
The supremum metric measures the largest coordinatewise discrepancy between
two bounded sequences.
\end{remark*}

\DefinitionalRoot

\begin{propositionbox}{Proposition}
\begin{proposition}[The Supremum Distance on Bounded Sequences is a Metric]
\label{prop:bounded-real-sequence-space-sup-metric}
The function \(d_\infty\) is a metric on \(\ell^\infty(\mathbb{R})\).

\smallskip
\noindent
\hyperref[prf:bounded-real-sequence-space-sup-metric]{\textit{Go to proof.}}
\end{proposition}
\end{propositionbox}

\begin{remark*}[Standard quantified statement]
\[
  \operatorname{IsMetric}(\ell^\infty(\mathbb{R}),d_\infty).
\]
\end{remark*}

\begin{remark*}[Predicate reading]
\[
\operatorname{IsMetric}(\ell^\infty(\mathbb{R}),d_\infty)
\]
means that the supremum distance is a metric on bounded real sequences.
\end{remark*}

\begin{remark*}[Interpretation]
Two bounded sequences are close in this metric when every coordinate is close
by the same tolerance.
\end{remark*}

\begin{dependencies}
\begin{itemize}
  \item \hyperref[def:metric-space]{Definition: Metric Space}
  \item \hyperref[def:bounded-real-sequence-space]{Definition: Supremum Metric on Bounded Real Sequences}
\end{itemize}
\end{dependencies}

\begin{definitionbox}{Definition --- Discrete Metric}
\begin{definition}[Discrete Metric]\label{def:discrete-metric}
Let \(X\) be a nonempty set.  The \emph{discrete metric} on \(X\) is the
function
\[
  d_{\mathrm{disc}}:X\times X\to\mathbb{R}_{\geq 0}
\]
defined by
\[
  d_{\mathrm{disc}}(x,y)
  =
  \begin{cases}
    0, & x=y,\\
    1, & x\neq y.
  \end{cases}
\]
\end{definition}
\end{definitionbox}

\begin{remark*}[Standard quantified statement]
\[
\begin{aligned}
&X\neq\varnothing,\qquad
d_{\mathrm{disc}}:X\times X\to\mathbb{R}_{\geq 0},\\
&\forall x,y\in X{:}\quad
d_{\mathrm{disc}}(x,y)=0 \Longleftrightarrow x=y,\\
&\forall x,y\in X{:}\quad
x\neq y \Longleftrightarrow d_{\mathrm{disc}}(x,y)=1.
\end{aligned}
\]
\end{remark*}

\begin{remark*}[Predicate reading]
\[
\operatorname{IsMetric}(X,d_{\mathrm{disc}})
\]
means that the discrete distance is a metric on the nonempty set \(X\).
\end{remark*}

\begin{remark*}[Interpretation]
The discrete metric treats equality as the only closeness relation.  Distinct
points are all placed at the same positive distance from each other.
\end{remark*}

\DefinitionalRoot

\begin{propositionbox}{Proposition}
\begin{proposition}[The Discrete Distance is a Metric]
\label{prop:discrete-metric-space}
For every nonempty set \(X\), the discrete distance
\(d_{\mathrm{disc}}\) is a metric on \(X\).

\smallskip
\noindent
\hyperref[prf:discrete-metric-space]{\textit{Go to proof.}}
\end{proposition}
\end{propositionbox}

\begin{remark*}[Standard quantified statement]
\[
X\neq\varnothing
\Longrightarrow
\operatorname{IsMetric}(X,d_{\mathrm{disc}}).
\]
\end{remark*}

\begin{remark*}[Predicate reading]
\[
\operatorname{IsMetric}(X,d_{\mathrm{disc}})
\]
means that the discrete distance is a metric on \(X\).
\end{remark*}

\begin{remark*}[Interpretation]
The only point at distance \(0\) from \(x\) is \(x\) itself.  Symmetry is built
into the two-case definition, and the triangle inequality is automatic because
the only possible distances are \(0\) and \(1\).
\end{remark*}

\begin{dependencies}
\begin{itemize}
  \item \hyperref[def:metric-space]{Definition: Metric Space}
  \item \hyperref[def:discrete-metric]{Definition: Discrete Metric}
\end{itemize}
\end{dependencies}

\begin{definitionbox}{Definition --- Singleton Subtype Metric}
\begin{definition}[Singleton Subtype Metric]\label{def:singleton-subtype-metric}
Let \((X,d)\) be a metric space and let \(a\in X\).  On the singleton
\(\{a\}\), the \emph{singleton subtype metric} is the restriction of \(d\) to
\(\{a\}\).  Equivalently, every pair of points in \(\{a\}\) has distance
\(0\).
\end{definition}
\end{definitionbox}

\begin{remark*}[Standard quantified statement]
\[
d_{\{a\}}:\{a\}\times\{a\}\to\mathbb{R}_{\geq 0},
\qquad
d_{\{a\}}(x,y)=d(x,y)=0.
\]
\end{remark*}

\begin{remark*}[Interpretation]
A singleton has only one point, so the only distance value forced by the
metric axioms is the distance from that point to itself.
\end{remark*}

\LeanFormalizes{def:singleton-subtype-metric}{lra-lean}{LRA.VolumeIV.MetricSpaces.Metrics}{SingletonMetric}{statement}

\DefinitionalRoot

\begin{definitionbox}{Definition --- Vacuous Metric on Empty Types}
\begin{definition}[Vacuous Metric on Empty Types]
\label{def:vacuous-metric-empty-types}
Let \(X\) be an empty type.  The \emph{vacuous metric} on \(X\) is the unique
function
\[
  d_{\varnothing}:X\times X\to\mathbb{R}_{\geq 0}.
\]
\end{definition}
\end{definitionbox}

\begin{remark*}[Standard quantified statement]
\[
X=\varnothing
\Longrightarrow
d_{\varnothing}:X\times X\to\mathbb{R}_{\geq 0}.
\]
All metric axioms hold vacuously because there are no points of \(X\).
\end{remark*}

\begin{remark*}[Interpretation]
An empty type has no pairs of points and therefore no nontrivial distance
values.  The metric axioms impose no additional conditions.
\end{remark*}

\LeanFormalizes{def:vacuous-metric-empty-types}{lra-lean}{LRA.VolumeIV.MetricSpaces.Metrics}{VacuousMetric}{statement}

\DefinitionalRoot

\begin{definitionbox}{Definition --- Empty Subtype Metric}
\begin{definition}[Empty Subtype Metric]\label{def:empty-subtype-metric}
Let \((X,d)\) be a metric space, and let
\[
  \varnothing_X:=\{x\in X : \text{false}\}.
\]
The \emph{empty subtype metric} is the restriction of \(d\) to
\(\varnothing_X\).
\end{definition}
\end{definitionbox}

\begin{remark*}[Standard quantified statement]
\[
d_{\varnothing_X}:
\varnothing_X\times\varnothing_X\to\mathbb{R}_{\geq 0},
\qquad
d_{\varnothing_X}(x,y)=d(x,y).
\]
\end{remark*}

\begin{remark*}[Interpretation]
The empty subtype metric is the inherited metric on an empty subset.  Since
the subset has no points, the inherited metric axioms are vacuous.
\end{remark*}

\LeanFormalizes{def:empty-subtype-metric}{lra-lean}{LRA.VolumeIV.MetricSpaces.Metrics}{EmptyMetric}{statement}

\DefinitionalRoot
